\setchapterpreamble[u]{\margintoc}
\chapter{T-SQL: Gatilhos}
\labch{cet615-gatilhos}

Este capítulo cobre a apresentação sobre gatilhos da terceira
aula~\cite{BaiocoCET615}. Com eles, a disciplina fecha a sequência de
objetos do banco vista desde a primeira aula: tabelas, índices,
restrições, visões, procedimentos e gatilhos.

\section{O que é um gatilho}

Um \textbf{gatilho}\index{gatilho} (\emph{trigger}) é um conjunto de
comandos SQL ligado a uma tabela e executado automaticamente pelo SGBD
quando os dados dela mudam por \texttt{insert}, \texttt{update} ou
\texttt{delete}. Pode ser visto como um procedimento armazenado sem
parâmetros, que ninguém chama: ele dispara sozinho a cada comando a que
está associado. Gatilhos servem para verificar ou modificar dados, e não
para devolver dados ao usuário com \texttt{select}.

O paralelo com o procedimento armazenado ajuda: os dois são conjuntos
de comandos SQL, mas o procedimento não pertence a uma tabela e é
chamado com \texttt{exec}; o gatilho fica preso a uma tabela e a um
comando, não recebe parâmetros e não pode ser chamado.

Usos convencionais citados na aula:

\begin{itemize}
	\item manter integridade que os recursos declarativos não cobrem
	(\texttt{unique}, \texttt{primary key}, \texttt{not null},
	\texttt{foreign key});
	\item propagar modificações entre tabelas relacionadas;
	\item manter valores derivados, como o total de uma nota fiscal
	calculado a partir dos itens.
\end{itemize}

\subsection{Bancos de dados ativos}

Os gatilhos também implementam \textbf{bancos de dados ativos}: bancos
que, além de guardar dados e executar transações, reagem a eventos. As
reações são regras ativas no modelo \textbf{evento--condição--ação}:

\begin{lstlisting}[language={},numbers=none]
on evento    -- a modificação que dispara a regra
if condição  -- teste opcional
then ação    -- o que fazer
\end{lstlisting}

Na história contada em aula\sidecite{BaiocoCET615Aulas}, os gatilhos
existiam nos gerenciadores antes do padrão SQL e nasceram para checar
restrições de integridade, numa época em que declarar \texttt{primary
key} só documentava a coluna. Quando os fabricantes passaram a checar as
restrições sozinhos, o gatilho ganhou outro uso: tornar o banco ativo. O
usuário que escolhe uma opção do menu não sabe que outros comandos rodam
por trás; o projetista e o programador sabem.

\textcite[seção 26.1]{Elmasri2011} descreve o mesmo modelo: o evento
costuma ser uma atualização no banco, mas pode ser um evento de tempo
(``todo dia às 5h30'') ou externo; a condição é opcional; e a ação pode
ser uma sequência de comandos SQL, uma transação ou um programa externo.

\section{Funcionamento no SQL Server}

Cada comando SQL é uma transação. Quando a tabela tem um gatilho para
aquele comando, a transação é estendida com o código do gatilho. Um
\texttt{delete} passa a valer como:

\begin{lstlisting}[language={},numbers=none]
begin transaction
   grava o log
   remove as linhas
   executa o gatilho de exclusão
commit ou rollback
\end{lstlisting}

Por isso um \texttt{rollback} dentro do gatilho desfaz também o comando
que o disparou, e um \texttt{rollback} em quem executou o comando desfaz
também o que o gatilho fez.

Três consequências que a professora destacou:

\begin{itemize}
	\item O gatilho não tem \texttt{begin transaction}, porque continua a
	transação do comando, nem \texttt{commit}, porque pode haver uma
	transação externa com outros comandos ainda por executar.
	\item Ela sempre põe um ponto de \texttt{rollback} no gatilho. Com
	mais de uma operação, testa a primeira antes de fazer a segunda, para
	não gastar tempo com operações inúteis. Esse \texttt{rollback} desfaz
	também a transação externa inteira.
	\item Gatilho não alivia bloqueio. Trocar uma transação explícita por
	um gatilho não deixa nada mais leve: a transação continua do mesmo
	tamanho. O ganho é de consistência, porque as ações rodam sempre que
	alguém executa o comando, sem depender de lembrar de abrir uma
	transação.
\end{itemize}

\begin{figure}[h]
\centering
\begin{tikzpicture}[node distance=4mm]
	\node[caixa, fill=white] (cmd) {insert into exemplar ...};
	\node[caixaazul, below=of cmd, minimum width=50mm] (exec) {grava log e altera a tabela};
	\node[caixaazul, below=of exec, minimum width=50mm] (rest) {verifica restrições (PK, FK, check)};
	\node[caixa, below=of rest, minimum width=50mm, fill=bdroxo, text=white] (trg) {executa o gatilho};
	\node[caixaazul, right=6mm of trg, minimum width=18mm] (ins) {inserted\\deleted};
	\node[caixa, below=of trg, minimum width=50mm] (fim) {commit ou rollback};
	\draw[seta] (cmd) -- (exec);
	\draw[seta] (exec) -- (rest);
	\draw[seta] (rest) -- (trg);
	\draw[seta] (trg) -- (fim);
	\draw[seta, dashed] (ins) -- (trg);
	\node[draw=bdcinza, dashed, fit=(exec)(fim), inner sep=2.5mm,
		label={[rotulo]left:uma transação}] {};
\end{tikzpicture}
\caption{Um gatilho \texttt{after} no SQL Server roda dentro da mesma
transação do comando, depois das restrições, uma vez por comando.}
\end{figure}

\subsection{As tabelas \texttt{inserted} e \texttt{deleted}}

Durante o gatilho, o SQL Server oferece duas tabelas lógicas, com a
mesma estrutura da tabela do gatilho, que podem ser consultadas com
\texttt{select}:

\begin{description}
	\item[\texttt{inserted}:] as linhas novas de um \texttt{insert}, ou os
	valores novos de um \texttt{update}.
	\item[\texttt{deleted}:] as linhas removidas por um \texttt{delete}, ou
	os valores antigos de um \texttt{update}.
\end{description}

Um \texttt{update} funciona como uma exclusão seguida de uma inclusão: o
valor antigo vai para \texttt{deleted} e o novo, para \texttt{inserted}.
Não existe uma tabela \texttt{updated}. Comparar as duas mostra o que
mudou; a diferença \texttt{i.quantidade - d.quantidade} serve para os
dois sentidos, porque fica negativa quando a quantidade diminui. No
Oracle, os equivalentes são \texttt{new} e \texttt{old}.

A função \texttt{update(coluna)} indica se a coluna aparece no
\texttt{set} do comando. Ela não compara valores: devolve verdadeiro
mesmo que o valor não mude. Serve para um gatilho de alteração mais
pontual, que só age quando a coluna que interessa foi tocada, já que
disparar à toa custa tempo e bloqueio em outra tabela.

Se o mesmo valor de \texttt{inserted} é usado várias vezes nas ações,
a professora sugere carregá-lo uma vez numa variável em vez de repetir o
\texttt{select}.

\section{Criar, alterar e excluir}

\begin{lstlisting}
create trigger nome
on tabela
for {[insert] [,] [update] [,] [delete]}
as
comandos

alter trigger nome on tabela for ... as ...  -- substitui a definição
exec sp_rename 'nome', 'novo_nome'
drop trigger nome1, nome2
\end{lstlisting}

O gatilho aparece no SSMS dentro da tabela, na pasta \emph{Gatilhos}, e
\texttt{sp\_helptext} mostra o código, como em visões e procedimentos.
\texttt{drop table} apaga os gatilhos da tabela junto; visões e
procedimentos precisam de \texttt{drop} próprio.

\subsection{\texttt{for}, \texttt{after}, \texttt{before} e
\texttt{instead of}}

O slide só mostra \texttt{for}, a forma antiga. \texttt{after} veio com
a padronização, e no SQL Server os dois são iguais: primeiro executa o
comando, depois o gatilho. O padrão e o Oracle têm também
\texttt{before}, que inverte a ordem; o SQL Server não tem. Ele tem
\texttt{instead of}, que \emph{substitui} o comando. Muita gente usa
\texttt{instead of} para simular \texttt{before}, repetindo o comando
original dentro do gatilho; a professora chama isso de adaptação, não do
propósito do recurso.

\subsection{Projetar antes de escrever}

``Não adianta saber sintaxe e não saber projetar um gatilho''. Antes de
escrever, defina quais operações, em quais tabelas, o disparam, e quais
ações ele executa. Duas orientações da aula:

\begin{itemize}
	\item Um gatilho por operação. A sintaxe aceita \texttt{insert,
	update, delete} no mesmo gatilho, mas isso só é comum em gatilhos de
	verificação; para regras ativas, as ações mudam de uma operação para
	outra (o \texttt{delete} subtrai o que o \texttt{insert} somou), e um
	gatilho único com \texttt{if} para escolher o código não compensa.
	\item Pode haver mais de um gatilho para a mesma operação. Ações
	independentes podem ficar separadas; ações em que a ordem importa
	ficam no mesmo gatilho, porque não se controla a ordem de execução
	entre gatilhos.
\end{itemize}

\subsection{Gatilho ou transação explícita?}

Tudo o que um gatilho faz pode ser feito por um procedimento com
transação; o contrário nem sempre. A professora deu um teste que, nas
palavras dela, não está nos slides: pense no procedimento que faria a
transação. Se os parâmetros da primeira operação bastam para executar
as outras, as outras podem virar gatilho da primeira, porque esses
valores estarão em \texttt{inserted} ou \texttt{deleted}.

\begin{itemize}
	\item \textbf{Dá para ser gatilho:} inserir um item de nota recebe
	número da nota, produto e quantidade. Com isso, o gatilho busca o
	preço e atualiza o total da nota, sem parâmetro extra. O mesmo vale
	para baixar o estoque ao vender um item.
	\item \textbf{Não dá:} cadastrar um atendente, que é uma pessoa
	genérica mais um registro específico. A segunda inclusão precisa de
	salário e comissão, que a primeira não tem; só o código se
	reaproveita.
\end{itemize}

O método sugerido para o projeto é listar tudo o que exige mais de uma
operação e decidir, caso a caso, o que fica em transação explícita e o
que vai para gatilho.

\section{Exemplos}

O exemplo da aula usa produto, nota fiscal e item de nota, e mantém o
valor total da nota com três gatilhos. Aqui o mesmo raciocínio aparece
no banco de biblioteca: a tabela \texttt{livro} ganha uma coluna
derivada \texttt{qtd\_exemplares}, mantida pelos gatilhos da tabela
\texttt{exemplar}.

\begin{lstlisting}
alter table livro add qtd_exemplares int not null default 0
\end{lstlisting}

\begin{lstlisting}[caption={Inclusão: soma os exemplares novos de cada livro.}]
create trigger exemplar_inclusao
on exemplar for insert
as
update livro
set qtd_exemplares = qtd_exemplares + i.qtd
from livro l
     inner join (select codlivro, count(*) as qtd
                 from inserted
                 group by codlivro) i
     on l.codlivro = i.codlivro
\end{lstlisting}

\begin{lstlisting}[caption={Exclusão: desconta os exemplares removidos.}]
create trigger exemplar_exclusao
on exemplar for delete
as
update livro
set qtd_exemplares = qtd_exemplares - d.qtd
from livro l
     inner join (select codlivro, count(*) as qtd
                 from deleted
                 group by codlivro) d
     on l.codlivro = d.codlivro
\end{lstlisting}

\begin{lstlisting}[caption={Alteração: só age se o exemplar mudou de livro.}]
create trigger exemplar_alteracao
on exemplar for update
as
if update(codlivro)
begin
   update livro set qtd_exemplares = qtd_exemplares - d.qtd
   from livro l inner join (select codlivro, count(*) as qtd
                            from deleted group by codlivro) d
        on l.codlivro = d.codlivro

   update livro set qtd_exemplares = qtd_exemplares + i.qtd
   from livro l inner join (select codlivro, count(*) as qtd
                            from inserted group by codlivro) i
        on l.codlivro = i.codlivro
end
\end{lstlisting}

O \texttt{update ... from} com junção é a forma do T-SQL de alterar uma
tabela usando dados de outra. O agrupamento por \texttt{codlivro} faz o
gatilho funcionar mesmo quando um único comando mexe em vários
exemplares; o porquê está no complemento a seguir.

Os gatilhos da aula seguem outro padrão: atualizam a nota com
\texttt{where numeronota = (select numeronota from inserted)} e, se
nenhuma linha foi afetada (\texttt{@@rowcount = 0}), executam
\texttt{rollback}. Na resolução dos exercícios, a professora observou
que ``nenhuma linha afetada'' quer dizer ``não encontrou'', e não
encontrar nem sempre é erro: quem decide abortar é a lógica do
gatilho.

O exercício de fixação traz mais um uso: quando uma fatura é paga
(\texttt{if update(dtpagamento)}), o gatilho a apaga de \texttt{fatura}
e a insere numa tabela \texttt{faturapaga} com \texttt{insert ...
select * from inserted}. O mesmo recurso alimenta tabelas históricas ou
desnormalizadas para BI, juntando \texttt{inserted} com tabelas reais.

\subsection{Gatilhos e chaves estrangeiras}

Os testes da aula mostram uma interação importante. Inserir um item
para uma nota que não existe falha pela chave estrangeira, antes de o
gatilho rodar. E o gatilho que apaga os itens de uma nota excluída nunca
chega a rodar: a chave estrangeira do item impede o \texttt{delete} da
nota antes. A solução dos slides é remover essa chave estrangeira e
passar a verificar, no gatilho de inclusão, se a nota existe
(\texttt{if not exists (...)}, com \texttt{print} e
\texttt{rollback}).

A professora fez isso de propósito, para mostrar como a cascata era
feita antes de existir \texttt{on delete cascade}, e avisou que não é
para repetir. Tirou só a chave estrangeira do item para a nota, mantendo
a do produto, e pegou o nome gerado da restrição na própria mensagem de
erro (conferido com \texttt{sp\_helpconstraint}). Sem a chave, um item
para a nota 0, que não existe, passou: uma inconsistência que ela chamou
de inaceitável e que o novo gatilho de inclusão passou a barrar. No fim,
o \texttt{delete} da nota apagou também os itens. Cascata e integridade
funcionaram, só que por dois gatilhos, e o gerenciador hoje entrega as
duas coisas prontas.

\section*{Complementos da documentação}

\subsection*{O gatilho dispara uma vez por comando}

Os slides dizem que o gatilho se repete para cada linha afetada. Em
aula, a professora ressalvou que isso depende do fabricante, lembrou que
o SQL Server não tem o \texttt{for each row} do Oracle e deixou para a
turma verificar o caso de várias linhas num comando. No SQL Server não é
por linha: o gatilho DML dispara \emph{uma vez por comando},
mesmo que o comando afete várias linhas ou nenhuma, e as tabelas
\texttt{inserted} e \texttt{deleted} trazem todas as linhas
afetadas\sidecite{MSCreateTrigger}. A documentação recomenda escrever
gatilhos preparados para várias linhas\sidecite{MSTriggerInsertedDeleted}.

A consequência prática: \texttt{where numeronota = (select numeronota
from inserted)} só funciona quando o comando insere um item por vez. Um
\texttt{insert ... select} com itens de duas notas faz a subconsulta
devolver dois valores, e o comando falha. A forma segura é juntar com
\texttt{inserted} e agrupar, como nos exemplos acima.

\textcite[seção 26.1]{Elmasri2011} explica a origem da confusão: no
padrão SQL e no Oracle, o gatilho pode ser declarado \texttt{for each
row} (dispara por linha) ou não (dispara por comando). O SQL Server só
tem a segunda forma.

\subsection*{Quando as linhas já estão na tabela}

Os slides dizem que \texttt{inserted} contém as linhas que serão
inseridas, mas ainda não foram. No SQL Server, \texttt{for} e
\texttt{after} são sinônimos: o gatilho roda \emph{depois} que o comando
foi executado com sucesso, inclusive as verificações de restrições e as
ações em cascata\sidecite{MSCreateTrigger}. Quando o gatilho roda, as
linhas já estão na tabela, dentro da transação ainda aberta; é isso que
o \texttt{rollback} desfaz. É também por isso que a chave estrangeira
barra o comando antes do gatilho no teste da aula.

\subsection*{Alternativas para manter a chave estrangeira}

Remover a chave estrangeira para o gatilho de exclusão funcionar troca
uma restrição declarativa por código. Duas alternativas mantêm a
integridade referencial:

\begin{itemize}
	\item declarar a chave estrangeira com \texttt{on delete cascade}
	(\refch{cet615-ddl}): o próprio SGBD apaga os itens, sem gatilho;
	\item usar um gatilho \texttt{instead of delete} na nota, que roda
	\emph{no lugar} do comando, antes das restrições: ele apaga os itens
	e depois a nota. A documentação não permite esse tipo de gatilho numa
	tabela que tenha \texttt{on delete cascade}\sidecite{MSCreateTrigger}.
\end{itemize}

Os gatilhos \texttt{instead of} também são a forma de tornar atualizável
uma visão que não seria (\refch{cet615-visoes}).

\subsection*{\texttt{commit} dentro do gatilho}

Em aula, a professora lembrou, como recordação incerta, que o SQL Server
ignora um \texttt{commit} no gatilho e só avisa. O teste mostra outra
coisa. Num gatilho de inclusão com \texttt{commit}, o SQL Server 2022
confirmou a inclusão e o que o gatilho fez, e em seguida deu o erro 3609
(\emph{the transaction ended in the trigger; the batch has been
aborted}): o resto do lote não rodou, e nenhuma transação externa
poderia mais desfazer aquelas linhas. É mais um motivo para a regra da
aula: gatilho não tem \texttt{commit}.

\subsection*{Gatilhos encadeados e recursivos}

Um aluno perguntou se um \texttt{update} dentro de um gatilho pode
disparar outro gatilho. Pode, se as opções do servidor permitirem. A
professora separou dois casos: o encadeamento, que vem ligado e em si não
causa problema; e a recursão, quando a cadeia volta a disparar um
gatilho anterior, que vem bloqueada porque, sem condição de parada, vira
laço infinito.

Os padrões batem com o SQL Server 2022: a opção de servidor
\texttt{nested triggers} vem ligada, e a opção de banco
\texttt{recursive\_triggers}, desligada\sidecite{MSCreateTrigger}. O
detalhe é o alcance: \texttt{recursive\_triggers} desligada só impede a
recursão \emph{direta}, o gatilho de T1 alterando T1. A indireta (T1
dispara T2, que altera T1) só é barrada desligando \texttt{nested
triggers}, o que barra todo encadeamento. Com ela ligada, um laço
indireto para no limite de 32 níveis, com erro.

\subsection*{Outros pontos}

\begin{itemize}
	\item Gatilhos que devolvem resultados com \texttt{select} estão
	marcados para remoção em versão futura.
	\item \texttt{truncate table} não dispara gatilho de exclusão.
	\item Sobre a ordem entre gatilhos da mesma operação: a aula diz que
	não há controle. Há um pouco: \texttt{sp\_settriggerorder} define o
	primeiro e o último gatilho \texttt{after}; os do meio rodam em
	ordem indefinida. A recomendação da aula, juntar no mesmo gatilho o
	que depende de ordem, continua a mais segura.
	\item Além dos gatilhos de dados, o SQL Server tem gatilhos de DDL,
	disparados por \texttt{create}, \texttt{alter}, \texttt{drop},
	\texttt{grant} e outros, no banco (\texttt{on database}) ou no
	servidor (\texttt{on all server}), e gatilhos de logon. A aula não
	tratou deles\sidecite{MSCreateTrigger}.
	\item Um gatilho mantém os bloqueios da transação enquanto roda, então
	deve ser curto. A documentação sugere sair logo quando o comando não
	afetou linha nenhuma: \texttt{if (rowcount\_big() = 0) return}.
\end{itemize}
